\section{Discussion}
\label{sec:discussion}

Our results identify a gap between retaining document information and using it
reliably. Incorrect D2L responses frequently match answers to other questions
about the same passage, while passage-token information remains detectable in
the compiler latent. Passage-specific fitting improves answers covered by its
supervision but degrades uncovered answers, suggesting that internalization
must preserve question--answer associations as well as content. Increasing
adapter rank alone did not resolve this problem at the tested fitting budget.

The distribution of composition damage depends on bank size. On Gemma, with eight
documents, interference accumulates across depth, and applying the full bank
in the initial layers provides little advantage over withholding document
contributions there. With sixteen documents, much of the damage has already
occurred by layer 8. Although this constrains late interventions, the
source-only oracle at $K=8$ retains substantial accuracy at intermediate depth,
so timing alone cannot account for the trained readers' failures.
The results identify a limitation of the tested composition and reader
configurations, while leaving open whether more effective selectivity can
be learned during compilation or reading.

Although compiling consolidated memory largely preserves its benefit on
PersonaMem, the resulting token savings are modest because episodic memories
remain text, and the serving measurements show that lower analytical compute
does not consistently translate into higher throughput. These limited gains
must be weighed against compilation cost and the attribution and interference
failures observed in factual QA, leaving little evidence that D2L offers a
compelling alternative to text-based memory in the production settings studied
here.

\subsection{Limitations and threats to validity}
We test only two models and one type of trained hypernetwork. It is unclear
whether other hypernetwork architectures or training methods would produce
the same failures.

\subsection{Future work: efficient online and continual learning}
At MidBrain, we want to turn semantic and procedural memories into internalized
behavior, so a model does not have to spend as much compute reading the same
memories in context on every request. D2L provides one way to do this, but the
failures studied here motivate continuing with other continual-learning
approaches.

We are particularly interested in learning from interactions in which a memory
is actually used. Such interactions provide both the request and an example
of how the memory should affect the response, offering supervision beyond the
memory text alone. For procedural memory, the objective is to learn how to
carry out a preference or instruction and recognize when it applies. Learning
across several interactions may help preserve this relationship while reducing
the need to repeatedly supply the instruction as text.

We believe online-learning may provide an interesting avenue of research.
Learning personalized behaviors at inference time could reduce the need for
separate training infrastructure. In principle, much of the information needed
for learning is already available during inference; if the necessary
activations are retained, what remains is largely a backward pass and parameter
update. However, online-learning also introduces challenges, including
catastrophic forgetting and difficult convergence properties.

Over time, these updates must accommodate new information without disrupting
behavior that remains valid. Our interest is therefore in bounded per-user
state that can accumulate compatible memories, adapt when preferences change,
and preserve the applicability of earlier instructions. We want to evaluate
this on sequences of interactions, measuring whether learning improves later
responses, how well previously acquired behavior is retained, and how much
context and serving computation can be avoided. This would connect the
immediate cost of an update to its value over the subsequent lifetime of the
memory.
